\scanpage{208}
\[
\unsure{G|\mathscr G}_k,\ \unsure{H|\mathscr H}_k\subset\Omega_k
\]
\[
G_k\cap H_k=\{u_\lambda\mid\lambda\in S_k,\ \lambda T=T\ \text{i.e.}\ (\lambda\{T\}-1)T=0\}
\]
Si on pose $\displaystyle\lambda\{T\}=\sum_{i\geq0}\lambda_iT^{(i)}$, alors la condition sur $\lambda$ envisagée est
\[
\boxed{\lambda_0=1}\quad\deleted{\unsure{\lambda_0=1},}\qquad\boxed{(i+1)\lambda_i=0\qquad\forall i\geq1}
\]
\deleted{ce qui implique en particulier $\lambda_0=1$.}

\deleted{\unsure{En car.} 0, cela implique \unsure{plus que}} Si $k$ sans torsion, cela implique $\lambda_i=0$ pour $i\geq1$ i.e. $\lambda=1$, i.e. $u_\lambda=\mathrm{id}$, d'où $G_k\cap H_k$\inserted{$=(\mathrm{id})$} dans ce cas\unsure{,|, si} --- et dans ce cas seulement. \deleted{\unsure{\}}}
[Dans ce cas \inserted{$k$ de car. 0 i.e. une $\mathbb Q$-alg.,}, $G_k$ est aussi le groupe \inserted{$G'_k$} des tous les endom. de $S_k=k\{T\}$ comme $k$-alg. augmentée \deleted{\unsure{p.d.}} et par suite \inserted{l'hom.} \deleted{\unsure{on a alors} $\Omega$}
\diagram{\begin{tikzcd}[ampersand replacement=\&,column sep=small]
(\unsure{1|0}\arrow[r]\&H'_k\arrow[r]\&{})\ \Omega_k\arrow[r,dash]\&G'_k=\operatorname{End}_{\substack{k\text{-alg}\\\deleted{\text{augm}}}}(\underset{\substack{\Vert\\S_k}}{\underbrace{k\{T\}}})\\
\&\&{}\&G_k\arrow[u,"(\unsure{??|\simeq|22})"']\arrow[ul,bend left=48]
\end{tikzcd}}
T
donne $\Omega_k$ comme produit $\frac12$ direct de $G_k$\inserted{$\simeq G'_k$} et de $H'_k$\ldots]

\deleted{Dans le cas général, on peut dire que}
\diagram{\begin{tikzcd}[ampersand replacement=\&,column sep=small]
1\arrow[r,dash]\&H_k\arrow[r,dash]\&\deleted{\unsure{G_k}}\Omega_k\arrow[r]\&\deleted{\unsure{\Omega'_k}}G'_k=\operatorname{End}_{k\text{-alg augm.}}(k\{T\})\\
\&\&{}\&G_k\arrow[u]\arrow[ul,bend left=48,dash,"{\unsure{\phantom{\rightarrow}|\rightarrow}}"']\\
\&\&{}\&G_k\cap H_k\simeq\{u_\lambda\mid\lambda=1+\sum_{i\geq1}\lambda_i T^{(i)},\ (i+1)\lambda_i=0\}\arrow[u]\\
\&\&{}\&{}\arrow[u]
\end{tikzcd}}
% archive sheet number: 207
